\documentclass[11pt]{article}
% cmap + T1 + a Type 1 text font so fi/fl ligatures survive copy-paste and
% full-text indexing. Without these, "five" extracts as " ve".
\usepackage{cmap}
\usepackage[T1]{fontenc}
\usepackage{lmodern}
\usepackage[utf8]{inputenc}
\usepackage[margin=1in]{geometry}
\usepackage{graphicx}
\usepackage{booktabs}
\usepackage{amsmath}
\usepackage[hidelinks]{hyperref}
\usepackage{xurl}
\usepackage{xcolor}
\usepackage{enumitem}
\usepackage{caption}
\usepackage{subcaption}
\usepackage{float}
\usepackage{listings}
\usepackage{natbib}
\bibliographystyle{unsrtnat}

\setlist{nosep,leftmargin=*}
\setlength{\parskip}{0.4em plus 0.1em minus 0.1em}

\lstset{
  basicstyle=\small\ttfamily,
  breaklines=true,
  frame=single,
  backgroundcolor=\color{gray!5},
  rulecolor=\color{gray!40},
  columns=fullflexible,
  keepspaces=true,
}

\definecolor{realgreen}{HTML}{2E7D32}
\definecolor{synthblue}{HTML}{1565C0}

\title{\textbf{IVF-Bench: A Rubric-Based Standard for Evaluating\\Vision-Language Models on IVF Clinical Reasoning}}
\author{
  Andrew G. A. Correa\\Monostate\\\texttt{andrew@monostate.ai}
  \and
  Brittany Yoon\\The Fertility Plan\\\texttt{b@thefertilityplan.com}
}
\date{August 2026}

\begin{document}
\maketitle

\begin{abstract}
Embryo selection in {IVF} remains a morphology-driven judgment, and the deep
learning systems built to standardize it operate on images alone, ranking
blastocysts without the patient context that clinicians actually weigh; the
largest randomized trial of that approach did not demonstrate noninferiority to
trained embryologists. We present IVF-Bench, a rubric-based benchmark for
evaluating whether vision-language models can perform the full assessment: 753
day-5 blastocyst cases from the public Kromp dataset, each pairing a real embryo
image and expert Gardner annotation with real cycle data and outcomes, together
with patient-history fields sampled from published population distributions and
explicitly labeled as generated, since no public dataset links such context to
embryo images. Models produce open-ended assessments graded on five clinical
rubrics by an {LLM} judge, and their stated implantation probabilities are scored
against recorded outcomes. Across seven frontier and open-weight systems,
clinical integration is uniformly strong while outcome discrimination is not,
with {AUROC} never exceeding 0.562 and every system scoring a worse Brier than a
constant set to the cohort pregnancy rate. Post-training Qwen 3.5-9B with {ORPO}
on 550 benchmark-derived preference pairs improves all five rubrics and places it
second of eight, ahead of Claude Opus 4.6 (4.11 against 4.01 held out,
$p{=}0.007$) at roughly one twentieth of the inference cost, indicating that the
reasoning component of embryo assessment does not require frontier-scale models,
while the prediction ceiling reflects the absence of context-linked outcome data
rather than a limitation of the models themselves. We release the benchmark, the
trained model, and all 5,194 responses with judge transcripts.
\end{abstract}

% ============================================================
\section{Introduction}

Few decisions in IVF carry more weight than the choice of which embryo to
transfer, and in most clinics that choice still rests on visual assessment: an
embryologist examines the blastocyst on day 5 and assigns a Gardner grade
capturing the degree of expansion, the organization of the inner cell mass, and
the cohesion of the trophectoderm \citep{gardner1999}. The grade is clinically
useful, but it is also subjective, and when embryologists at different centres
are shown the same day-5 embryos and asked to select one for transfer, their
agreement falls well short of what the confidence of the decision would
suggest \citep{storr2017}.

That variability has motivated a decade of machine learning work, nearly all of
it framed as a scoring problem in which an image or a time-lapse sequence is
mapped to a scalar and embryos are ranked by that scalar. The largest randomized
evaluation of this paradigm, a double-blind trial across 14 clinics and 1,066
patients, compared a deep learning ranker against trained embryologists and did
not demonstrate noninferiority \citep{illingworth2024}, which suggests that
further refinement of image-only ranking is unlikely to be the most productive
direction.

We argue that the limitation lies in the framing rather than in the
architecture. The embryo image captures only part of what determines whether a
transfer succeeds; the remainder is patient-level, comprising ovarian reserve,
endometrial receptivity, body mass index and lifestyle, the history of previous
cycles, and the andrology profile. In clinical practice most of this information
resides with the treating physician and rarely reaches the embryologist at the
microscope, still less the imaging model, so a system operating on pixels alone
is not so much failing at prediction as being asked to predict without the
covariates that govern the outcome.

Vision-language models are, in principle, capable of consuming both modalities,
but the claim has never been testable, because the data required to test it are
not publicly available: embryo images are sensitive clinical records,
outcome-linked patient context is more sensitive still, and the regulatory and
privacy requirements for assembling the two together are stringent enough that
essentially every substantial dataset in this field remains private. This, rather
than model capacity, is the binding constraint, and it motivates the first
contribution of this work, which is an instrument for measuring whether a model
can use patient context at all.

\textbf{IVF-Bench} is an open benchmark of 753 day-5 blastocyst cases. Each case
pairs a real embryo image and its expert Gardner annotation with real cycle and
outcome data from the same patient, plus the patient-history fields that no
public dataset carries. Those we generate from published population
distributions, deterministically, and label every one of them as generated. A
model is asked to describe the embryo, integrate the patient factors, estimate an
implantation probability, and recommend next steps. Five rubrics score the answer
from 1 to 5, and the stated probability is checked against the real outcome.

Generating that context is a deliberate methodological choice and we defend it in
Section~\ref{sec:synthetic}. It is what makes a case-level reasoning task
constructible from public data at all. It also has a consequence we state up
front: our generated fields are sampled independently of the recorded outcomes,
so they carry no information about who conceived. Any outcome-prediction result
in this paper is bounded by that before a model is involved.

We evaluate seven frontier and open-weight models, then post-train a 9B open
model on preference pairs drawn from the benchmark. What we found:

\begin{enumerate}
\item Targeted post-training substitutes for scale on this task: 550 preference
      pairs are sufficient to move a 9B open model past Claude Opus 4.6 into
      second position of eight at roughly one twentieth of the inference cost,
      indicating that the reasoning component of embryo assessment is accessible
      to models far below frontier scale.
\item The improvement is distributed across all five rubrics rather than
      concentrated in one, which distinguishes a procedure that improves the
      assessment as a whole from one that optimizes a single reported quantity.
\item Outcome discrimination remains weak for every system evaluated, and we
      quantify the deficit against a baseline that ignores its inputs entirely,
      attributing the ceiling to the available data rather than to model
      capacity, for the reason given above.
\item Morphology grounding is the weakest rubric throughout, though our judge
      does not observe the image, so we treat those scores as an upper bound on
      image comprehension rather than a measurement of it.
\end{enumerate}

The practical implication is that progress in this area depends less on model
development than on the construction of a public dataset linking embryo images to
measured patient context and recorded outcomes, and IVF-Bench provides the
instrument with which the value of such a dataset can be quantified.

% ============================================================
\section{Related work}

\textbf{Embryo selection models.} Most published systems predict a label from
pixels. They grade blastocysts automatically, rank embryos within a cohort, or
estimate implantation from time-lapse video. Commercial products in clinics today
follow the same pattern and are positioned as adjunctive scoring tools rather
than autonomous decision makers. The randomised evidence for the strongest of
them is the trial discussed above \citep{illingworth2024}. None of this work
evaluates open-ended clinical reasoning, because none of it produces language.

\textbf{Language about embryos.} The group that released the dataset we build on
has recently moved in this direction, publishing an expert-annotated set of
embryo images paired with natural-language morphological descriptions
\citep{neu2026dataset} and a fine-tuned PaliGemma-2 model that generates such
descriptions \citep{neu2026invitrovision}. That work and ours are complementary
and different in an important way. They evaluate whether a model can describe an
embryo. We evaluate whether a model can reason about a patient who has that
embryo, which requires the cycle data, the history, a probability, and a
recommendation. Their descriptions are the ground truth we most wish we had for
scoring morphology, and we say so in Section~\ref{sec:limitations}.

\textbf{Rubric-graded evaluation of medical language models.} Scoring open-ended
clinical answers against written criteria, with a language model doing the
scoring, was put on a solid footing by HealthBench, which had 262 physicians
write 48,562 rubric criteria over 5,000 conversations
\citep{arora2025healthbench}. Our five rubrics are far smaller and narrower, and
are written for one decision rather than for medicine in general. The
methodological debt is direct and we follow the same basic recipe: grade against
criteria, not against a gold answer.

\textbf{Preference optimisation.} We post-train with ORPO
\citep{hong2024orpo}, which folds the preference objective into the supervised
loss and removes the need for a separate reference model, in the family of
methods opened up by DPO \citep{rafailov2023dpo}.

% ============================================================
\section{The benchmark}

\subsection{Where the data comes from}

Open embryo data is rare, because embryo images are sensitive clinical records
and most published systems are trained on private single-clinic collections. We
build on the Kromp blastocyst dataset \citep{kromp2023}, which is public under CC
BY 4.0 and links day-5 embryo images to expert Gardner grades, selected cycle
variables, and pregnancy outcomes. Of the 837 patients in the release, 754 have
complete clinical records, and one of those references an image that is missing
from the distributed archive, which leaves 753 usable cases.

Preparing the source data took more work than we expected, and we document it
because anyone reusing this dataset will hit the same problems. The clinical
annotation file is semicolon-separated, Latin-1 encoded, and was clearly opened in
a German-locale spreadsheet at some point: 500 AMH values and 144 endometrial
thickness values had been silently converted into dates, so that a value of 1.26
appears in the file as \texttt{J\"an.26}. We map the month abbreviations back to
numbers to recover them. Oocyte counts sometimes appear in additive form such as
\texttt{7+1}, which we sum. Missing values stay missing and are shown to models as
"not available" rather than being imputed. One record has a live birth recorded
without the clinical pregnancy that must precede it; we keep it as it is, because
silently fixing a source dataset is worse than reporting the inconsistency.

\subsection{What a case contains}

Each case combines real fields from Kromp with generated patient context. The
real fields are the embryo image, the Gardner grades, patient age, AMH,
endometrial thickness, oocytes retrieved, mature oocytes, transfer day, and the
three pregnancy outcomes. The generated fields are BMI, basal FSH, diagnosis,
stimulation protocol, previous cycles and their outcomes, partner age, semen
parameters, and a lifestyle summary.

\subsection{Why the patient context is generated, and what that costs}
\label{sec:synthetic}

An embryo image plus an age is not a case. A clinician reasoning about a transfer
uses a history, and a benchmark that omits the history cannot tell whether a
model uses one. But no public dataset carries BMI, lifestyle, stimulation
protocol, or cycle history alongside embryo images and outcomes. That data exists
in clinics, mostly in physicians' notes and often in physicians' heads, and it
does not travel: it rarely reaches the embryologist at the microscope, and it has
never reached a public research corpus, because assembling images plus
outcome-linked patient records is exactly the combination that regulatory and
privacy review is built to slow down.

We therefore generate these fields, sampling each from a published population
distribution, with every parameter and its source given in
Appendix~\ref{app:distributions}. Generation is deterministic per case, seeded
from a hash of the case identifier, so that cases regenerate exactly and the
addition of new cases does not perturb existing ones, and each case record
enumerates its generated fields explicitly, so that no result can depend on
generated values being mistaken for measured ones.

This procedure yields a case-level reasoning task that can be constructed and
reproduced entirely from public sources, at two costs that bear directly on the
interpretation of our results. First, the generated fields are sampled
independently of the recorded outcomes, so that although their marginal
distributions are realistic they carry no information about which patients
conceived, and no model can extract predictive signal from them because none is
present; the outcome metrics in Section~\ref{sec:results} are consequently
bounded by the data before any model is involved, and should be read as a floor
that quantifies what is missing rather than as a verdict on what vision-language
models can achieve given measured context. Second, the fields are sampled
independently of one another, so genuine clinical correlations, such as that
between polycystic ovary syndrome and body mass index, are absent, and a model
that has internalized those correlations receives no credit for them here.
Neither limitation can be remedied by a more sophisticated generator; both are
remedied by a real dataset, which is precisely the argument of this paper.

\subsection{Splits}

The 753 cases are split deterministically into 550 test, 100 validation, and 103
held out. The test split carries the main leaderboard. The validation split is
reserved for judge and threshold work. The held-out split was untouched through
all training, hyperparameter search, and judge calibration, and exists to answer
one question: do the test-split rankings survive on cases nobody optimised
against?

One property of the splits is worth stating plainly, because it is a defect. They
were drawn per image rather than per patient. Twenty-two patients contributed two
embryos each, and for six of those patients the two embryos landed in different
splits. Four of the 103 held-out cases therefore share a patient with a test-split
case. We measure the effect in Section~\ref{sec:robustness} and it changes no
ranking, but a future version should split on patient.

\subsection{The prompt}

Every model sees the embryo image and a structured case description containing
the Gardner grade with plain-language descriptors, the laboratory fields, and the
patient profile. It is asked for five things: a morphology description, an
integration of the patient factors, the risks and favourable features, a numeric
implantation probability, and recommendations for this transfer and for future
cycles. The full template is in Appendix~\ref{app:prompt}.

\subsection{The five rubrics}

A judge model reads the case and the response and returns a 1 to 5 score with a
short justification for each of five rubrics.

\begin{itemize}
\item \textbf{Morphology grounding.} Does the description of the embryo match the
      image and the Gardner annotation?
\item \textbf{Clinical integration.} Does the response connect the patient factors
      to this embryo and this prognosis, rather than listing them?
\item \textbf{Reasoning coherence.} Is the argument internally consistent, with
      uncertainty where uncertainty belongs?
\item \textbf{Guideline alignment.} Is the advice consistent with evidence-based
      reproductive medicine?
\item \textbf{Recommendation usefulness.} Could a clinician act on this, or is it
      generic?
\end{itemize}

Full rubric text and the 1 to 5 anchors for each are in
Appendix~\ref{app:rubrics}. The judge also extracts the numeric probability from
the response, which we compare against the recorded outcome.

% ============================================================
\section{Evaluation setup}
\label{sec:setup}


We evaluate seven models through their public APIs plus our own post-trained
model served locally, and we grade every response with the same judge.

\begin{table}[H]
\centering
\small
\setlength{\tabcolsep}{5pt}
\begin{tabular}{llll}
\toprule
\textbf{Model} & \textbf{Served by} & \textbf{Size} & \textbf{Reasoning} \\
\midrule
GPT-5.4 (2026-03-05) & OpenAI & not disclosed & effort: high \\
Claude Opus 4.6 & Bedrock & not disclosed & 8K thinking budget \\
Claude Sonnet 4.6 & Bedrock & not disclosed & 8K thinking budget \\
Qwen 3.5-397B-A17B & OpenRouter & 397B, 17B active & native \\
Kimi K2.5 & OpenRouter & 1T, 32B active & native \\
Gemini 2.5 Flash & OpenRouter & not disclosed & native \\
Qwen 3.5-9B & OpenRouter & 9B & native \\
\midrule
\textbf{IVF-Bench-Qwen9B-ORPO} (ours) & local vLLM, 2$\times$H100 & 9B & native \\
\bottomrule
\end{tabular}
\caption{The eight systems evaluated. Baselines run at temperature 0 with
reasoning enabled. Our model uses the sampling settings recommended on the Qwen
3.5 card (temperature 1.0, top-p 0.95, top-k 20, repetition penalty 1.05),
because greedy decoding traps it in a repetition loop.}
\label{tab:models}
\end{table}

\textbf{The judge.} GPT-5.4 grades every response against the five rubrics in a
single call and also extracts the numeric probability. It sees the case text and
the response, but not the embryo image. This is a deliberate limitation and we
are explicit about what it costs us: the morphology rubric measures whether a
description is consistent with the Gardner annotation the model was handed, not
whether the model read the image correctly. A judge that cannot see the image
cannot catch a fluent description of an embryo that is not there.

\textbf{Two evaluation defects and their correction.} We report both because
each is easy to reproduce inadvertently and neither is visible in the resulting
scores.

The first concerns failed judge calls. When a call errored, the pipeline
persisted a score record containing five zeros and continued, and because
subsequent runs skip cases for which a score record already exists, those
failures became permanent. In our initial pass this silently invalidated 231 of
the test-split scores, 211 of which belonged to a single model, so that the
model's reported row was computed from 309 of 550 cases while the accompanying
table indicated 520. The judge now declines to persist a failed score, leaving
the case pending for a subsequent run.

The second concerns how reasoning traces are delivered. Our post-trained model
was served through vLLM without a reasoning parser, and because Qwen chat
templates pre-fill \texttt{<think>} in the generation prompt, the model emits its
reasoning, a closing \texttt{</think>}, and then its answer within a single
field, whereas every hosted API returns reasoning in a separate field. Our model
was consequently graded on approximately 4,000 characters of unfiltered reasoning
prepended to its answer while all other systems were graded on the answer alone;
correcting this raised its test-split score from 3.80 to 4.17 and moved it from
third position to second. The general point is that comparisons between locally
served reasoning models and hosted APIs require verifying the contents of the
string being graded.

All results reported below follow both corrections: every score in the study is
produced by GPT-5.4, and no score record contains a judge failure.


% ============================================================
\section{Results}
\label{sec:results}


\subsection{Reading the two outcome metrics}

\textbf{Brier score} (lower is better) is the mean squared error between the
stated probability and what happened. It rewards being honest about the odds.
Always saying 50\% scores 0.25, which is the comparison most papers use and the
wrong one: the sharper test is a constant set to the cohort's own pregnancy rate,
and we report against that in Table~\ref{tab:baserate}. \textbf{AUROC} (higher is
better, 0.5 is chance) is the probability that a randomly chosen success is
ranked above a randomly chosen failure. It rewards telling embryos apart.
Counselling a patient needs the first. Choosing which embryo to transfer needs
the second, and a model can be respectable on one while useless on the other.

\subsection{The leaderboard}

\begin{table}[H]
\centering
\small
\setlength{\tabcolsep}{4pt}
\begin{tabular}{clcccccccc}
\toprule
\textbf{\#} & \textbf{Model} & \textbf{Overall} & \textbf{Morph} & \textbf{Clin} &
\textbf{Reas} & \textbf{Guide} & \textbf{Rec} & \textbf{Brier}$\downarrow$ &
\textbf{AUROC}$\uparrow$ \\
\midrule
1 & GPT-5.4 & \textbf{4.54} & \textbf{3.96} & \textbf{5.00} & \textbf{4.66} & \textbf{4.15} & \textbf{4.92} & 0.245 & 0.553 \\
2 & \textbf{Qwen9B-ORPO (ours)} & 4.17 & 3.60 & 4.83 & 4.00 & 3.96 & 4.44 & 0.249 & 0.531 \\
3 & Opus 4.6 & 4.01 & 3.52 & 4.97 & 3.96 & 3.20 & 4.38 & 0.245 & 0.541 \\
4 & Qwen 397B & 3.75 & 3.33 & 4.27 & 3.84 & 3.31 & 3.98 & 0.260 & 0.538 \\
5 & Kimi K2.5 & 3.63 & 3.39 & 4.38 & 3.66 & 3.00 & 3.72 & \textbf{0.242} & 0.552 \\
6 & Sonnet 4.6 & 3.58 & 2.92 & 4.48 & 3.64 & 2.79 & 4.06 & \textbf{0.241} & 0.544 \\
7 & Gemini 2.5 Flash & 3.40 & 3.34 & 3.88 & 3.58 & 3.53 & 2.66 & 0.269 & 0.537 \\
8 & Qwen 9B (base) & 3.35 & 3.08 & 3.93 & 3.38 & 2.81 & 3.56 & 0.255 & \textbf{0.554} \\
\bottomrule
\end{tabular}
\caption{Test split, 550 cases, GPT-5.4 judge, every model with reasoning
enabled. Qwen 9B base completed 520 of 550 cases; the other 30 ran out of tokens
inside a reasoning loop. Rubrics are morphology grounding, clinical integration,
reasoning coherence, guideline alignment, and recommendation usefulness.}
\label{tab:leaderboard}
\end{table}

\begin{table}[H]
\centering
\small
\setlength{\tabcolsep}{4pt}
\begin{tabular}{clcccccccc}
\toprule
\textbf{\#} & \textbf{Model} & \textbf{Overall} & \textbf{Morph} & \textbf{Clin} &
\textbf{Reas} & \textbf{Guide} & \textbf{Rec} & \textbf{Brier}$\downarrow$ &
\textbf{AUROC}$\uparrow$ \\
\midrule
1 & GPT-5.4 & \textbf{4.55} & \textbf{3.98} & \textbf{5.00} & \textbf{4.65} & \textbf{4.23} & \textbf{4.88} & 0.246 & \textbf{0.562} \\
2 & \textbf{Qwen9B-ORPO (ours)} & 4.11 & 3.56 & 4.76 & 3.93 & 3.93 & 4.37 & 0.245 & 0.547 \\
3 & Opus 4.6 & 4.01 & 3.60 & 4.92 & 3.98 & 3.22 & 4.31 & 0.247 & 0.551 \\
4 & Gemini 2.5 Flash & 3.83 & 3.34 & 4.40 & 3.90 & 3.54 & 3.98 & 0.289 & 0.504 \\
5 & Qwen 397B & 3.82 & 3.48 & 4.31 & 3.88 & 3.44 & 3.98 & 0.265 & 0.551 \\
6 & Kimi K2.5 & 3.73 & 3.48 & 4.42 & 3.80 & 3.04 & 3.90 & 0.254 & 0.518 \\
7 & Sonnet 4.6 & 3.65 & 3.07 & 4.51 & 3.73 & 2.89 & 4.06 & \textbf{0.242} & 0.534 \\
8 & Qwen 9B (base) & 3.31 & 2.99 & 3.91 & 3.18 & 2.85 & 3.64 & 0.285 & 0.450 \\
\bottomrule
\end{tabular}
\caption{Held-out split, 103 cases that were never used for training,
hyperparameter search, or judge calibration. All eight models answered all 103.}
\label{tab:heldout}
\end{table}

\subsection{Which gaps are real}

A leaderboard without error bars invites readers to take a 0.02 difference
seriously. We resample cases with replacement 10,000 times and, for each adjacent
pair of models, bootstrap the paired per-case difference on the cases both models
answered.

On the test split every adjacent gap is separable at the 95\% level, the
narrowest being Kimi K2.5 over Sonnet 4.6 at $+0.052$ (95\% CI $[+0.005,
+0.098]$, $p{=}0.031$). On the held-out split, with 103 cases instead of 550, two
gaps disappear: Gemini Flash over Qwen 397B ($+0.015$, CI $[-0.049, +0.082]$,
$p{=}0.66$) and Kimi K2.5 over Sonnet 4.6 ($+0.076$, CI $[-0.023, +0.169]$,
$p{=}0.12$). Those pairs should be read as ties. The two results we care about
most survive on both splits: GPT-5.4 over our model ($+0.44$ held out, CI
$[+0.38, +0.51]$) and our model over Opus 4.6 ($+0.10$ held out, CI $[+0.03,
+0.17]$, $p{=}0.007$).

\subsection{What the numbers say}

\textbf{1. Case explanation is uniformly strong.} Clinical integration is the
highest-scoring rubric for every system evaluated, with GPT-5.4 reaching a
maximal 5.00 on both splits; models reliably incorporate age, AMH, endometrial
thickness, and prior cycle history into a coherent account of the patient's
prognosis, such that if the task were confined to producing the counselling
narrative it would be close to solved.

\textbf{2. Outcome prediction sits at the floor, and we can say exactly how far
below useful that is.} AUROC spans 0.450 to 0.562 across the field, where 0.5 is
chance. The best number in the study, 0.562, means that given one embryo that led
to a clinical pregnancy and one that did not, the best model ranks them correctly
56\% of the time.

Brier scores look better than they are. They cluster between 0.241 and 0.289,
which beats the 0.25 you get by always saying 50\%, and it is tempting to read
that as decent calibration. It is not the right comparison. The clinical
pregnancy rate in this cohort is 0.350, so a constant predictor that says
"35\%" for every case, ignoring the embryo, the image, and the patient entirely,
scores 0.2274. Every model in the study loses to it.

\begin{table}[H]
\centering
\small
\begin{tabular}{lrr}
\toprule
\textbf{Predictor} & \textbf{Brier (test)} & \textbf{Brier (held out)} \\
\midrule
Constant, the cohort base rate of 0.35 & \textbf{0.2278} & \textbf{0.2274} \\
Constant, 0.50 & 0.2500 & 0.2500 \\
\midrule
GPT-5.4 & 0.2447 & 0.2465 \\
Qwen9B-ORPO (ours) & 0.2491 & 0.2447 \\
Qwen 9B (base) & 0.2552 & 0.2851 \\
\bottomrule
\end{tabular}
\caption{Brier score against a predictor that ignores its inputs. No model in
this study beats the cohort base rate on either split. Model rows are computed on
the cases where a probability could be extracted.}
\label{tab:baserate}
\end{table}

This is the number we most want a reader to take away, and it is not a verdict on
the models. Section~\ref{sec:synthetic} explains why. The patient-context fields
we added are sampled independently of outcome, so they contain no signal any
model could use, and what remains of the real data is age, AMH, endometrial
thickness, oocyte counts, and one morphology grade from a single day-5 frame.
Predicting a live birth from that is close to the limit of what the public data
supports, and no architecture fixes it. The inputs that would move this number
are embryo genetics, longitudinal development rather than one frame, uterine and
endometrial context, and the lifestyle and history data that today never leaves
the physician's notes. IVF-Bench is built so those fields can be swapped from
generated to measured as they become available, which is exactly the experiment
the field cannot currently run.

The same reasoning reframes the calibration result for our post-trained model in
Section~\ref{sec:orpo}. Its held-out Brier of 0.245 is second of the eight,
behind Sonnet 4.6 at 0.242, and all eight sit below a constant. What
post-training bought is a model whose stated probabilities moved back toward the
population rate, correcting the base model's overconfidence. On this data that is
the right behaviour, and it is not predictive skill.

\textbf{3. Morphology grounding is the weakest rubric throughout.} No system
exceeds 3.98 on morphology grounding on either split, whereas five of eight
exceed 4.30 on clinical integration on the test split, indicating that models
reason more effectively about the structured data supplied in text than about the
image itself. Because our judge does not observe the image, a system can score
well on this rubric by restating the Gardner grade it was given, so these values
constitute an upper bound on image comprehension rather than a measurement of
it.

\textbf{4. Cost and quality are only loosely coupled.} On the test split Opus
4.6 requires \$184.81 to process the 550 cases and scores 4.01, while our 9B
model requires \$7.98 and scores 4.17; Gemini 2.5 Flash, the least expensive
commercial system evaluated, finishes seventh on that split while rising to
fourth on held out, the largest between-split movement of any system.


% ============================================================
\section{Post-training a small model}
\label{sec:orpo}


The natural question raised by the gap between a 9B open model and a frontier
system is whether that gap reflects a capability ceiling or a deficit of
domain-specific training data, and we tested this directly.

\subsection{Building preference pairs from the benchmark}

For each of the 550 test-split cases we already had seven scored responses. We
took the highest-scoring one as \texttt{chosen} and the lowest-scoring one as
\texttt{rejected}, giving 550 pairs with a mean score gap of 1.49 points, split
500 for training and 50 for evaluation. We then fine-tuned Qwen 3.5-9B with ORPO
\citep{hong2024orpo}, LoRA rank 16 and alpha 32 applied to all linear layers
including the vision tower, three epochs on two H100s, with the learning rate and
the preference weight chosen by an eight-trial search.

Two properties of this dataset matter for interpreting the result. The pairs come
from the test split, so the trained model is in-domain on the 550-case
leaderboard and the held-out split is the honest measurement. And the chosen
responses are close to a monoculture: GPT-5.4 wrote 460 of the 500 training
pairs' chosen side, 92\%, which means the student is largely distilling one
teacher, and that teacher is also the judge.

\subsection{What it bought}

\begin{table}[H]
\centering
\small
\begin{tabular}{lrrr}
\toprule
\textbf{Rubric} & \textbf{Base 9B} & \textbf{After ORPO} & \textbf{Change} \\
\midrule
Overall & 3.31 & \textbf{4.11} & $+24.0\%$ \\
Morphology grounding & 2.99 & 3.56 & $+19.2\%$ \\
Clinical integration & 3.91 & 4.76 & $+21.6\%$ \\
Reasoning coherence & 3.18 & 3.93 & $+23.5\%$ \\
Guideline alignment & 2.85 & 3.93 & $+38.2\%$ \\
Recommendation usefulness & 3.64 & 4.37 & $+20.0\%$ \\
\midrule
Brier score $\downarrow$ & 0.285 & \textbf{0.245} & $-14.2\%$ \\
AUROC $\uparrow$ & 0.450 & 0.547 & $+21.7\%$ \\
\bottomrule
\end{tabular}
\caption{Base Qwen 3.5-9B against the same model after ORPO on 550 preference
pairs, measured on the 103-case held-out split. Every rubric improves, guideline
alignment most of all.}
\label{tab:orpo}
\end{table}

Five hundred and fifty preference pairs moved a 9B model from last position to
second, ahead of Claude Opus 4.6, and every rubric improved, which is the result
worth emphasizing: the procedure raises the assessment as a whole rather than
optimizing a single reported quantity, and it is reproducible by anyone with the
released dataset and two GPUs.

Its held-out Brier of 0.245 is second of the eight, behind Sonnet 4.6 at 0.242,
and should be read alongside Table~\ref{tab:baserate}. What the fine-tune fixed
is overconfidence: the base model's mean stated probability was 0.452 against a
true rate of 0.350, and the trained model's is 0.396. Whether the same recipe yields real predictive skill
when the context fields are measured rather than generated is the open question
this benchmark exists to make answerable.

The base model's AUROC is instructive in the opposite direction: it attains
0.554 on the test split, the highest of the eight, and collapses to 0.450 on held
out, below chance, which is the signature of a ranking that has overfitted to one
distribution and fails to survive a shift, whereas the post-trained model moves
from 0.531 to 0.547, suggesting that the fine-tune regularized its probability
estimates rather than sharpening them.

The cost structure is the consideration most relevant to deployment: processing
the 550 test cases requires \$184.81 with Opus 4.6 against \$7.98 with our model,
which scores higher, and at approximately 19 seconds and \$0.017 per case on
hardware already provisioned, a post-trained small model is the appropriate form
factor for this workload, provided that the requirement is explanation rather
than prediction.


% ============================================================
\section{What the benchmark itself gets wrong}
\label{sec:robustness}


A benchmark paper should attempt to falsify its own instrument, and we report
four such checks, each of which reuses the collected responses and therefore
incurs no additional cost.

\textbf{Do the five rubrics measure five distinct constructs?} Only partly: the
mean correlation between distinct rubrics is 0.49 and the first principal
component accounts for 60\% of the variance across 4,370 judgements, indicating a
substantial general quality factor. The structure is nonetheless not degenerate,
since recommendation usefulness correlates with clinical integration at 0.69 but
with guideline alignment at only 0.28, which justifies reporting the five
separately while cautioning against treating their unweighted mean as five
independent measurements.

\textbf{Does the judge reward verbosity?} It does, and unevenly: pooled across
systems the Spearman correlation between response length and overall score is
0.61, and per system it ranges from $-0.01$ for Opus 4.6 to $+0.76$ for Gemini
2.5 Flash, which is both the lowest-scoring commercial system and the one whose
score depends most heavily on output length. Future versions of this benchmark
should control for length explicitly, and readers comparing closely spaced scores
should consider whether one system simply produced more text.

\textbf{Does patient-level leakage affect the results?} It does not: twenty-two
patients contributed two embryos each, and retaining only one embryo per patient
removes 16 cases from the test split without changing any ranking or moving any
score by more than 0.01. The defect is real and should be addressed by splitting
on patient rather than on image, but it did not contaminate these results.

\textbf{Does judge self-preference inflate our result?} We cannot exclude it,
and this remains the principal open question in the paper: GPT-5.4 serves as the
judge, produced 92\% of the chosen responses from which our model was trained,
and occupies the top of the leaderboard, so a model optimized to resemble the
judge should be expected to satisfy it. Confirming the ordering with a judge from
a different model family is the most immediate extension of this work, and until
that is done our model's second position should be read as second according to
GPT-5.4."


% ============================================================
\section{Limitations}
\label{sec:limitations}


\begin{itemize}
\item \textbf{The judge cannot see the embryo.} Morphology grounding therefore
      measures consistency with the Gardner grade in the prompt, not fidelity to
      the image. The natural fix is the expert-written morphological descriptions
      released by \citet{neu2026dataset}, which would let a future version score
      image reading against something real.
\item \textbf{One judge, and it is also a contestant.} GPT-5.4 grades a
      leaderboard it appears on, and it wrote 92\% of the responses our model was
      trained to imitate. Scoring a subset with a judge from a different family
      would settle how much that inflates our second place. It is unfinished
      work, not a design choice.
\item \textbf{The prompt asks for implantation probability and we score against
      clinical pregnancy.} Both endpoints are computed in the code; the
      leaderboard reports clinical pregnancy. Treat the outcome metrics as a
      stress test rather than a matched forecast.
\item \textbf{Part of the patient context is generated}, for the reasons and with
      the consequences set out in Section~\ref{sec:synthetic}. The two that bind
      hardest: those fields are independent of outcome, so they cap what any
      outcome metric here can show, and they are independent of each other, so
      real clinical correlations such as PCOS with BMI are absent.
\item \textbf{One clinic, one scanner, one population.} Everything here comes from
      the Kromp release. Transfer to other labs is untested.
\item \textbf{Splits are per image rather than per patient.} Quantified above;
      fix in the next version.
\item \textbf{Small absolute sample.} 550 test cases and 103 held out is enough to
      separate most models and not enough to separate close ones, which is why
      every comparison in this paper carries an interval.
\end{itemize}


% ============================================================
\section{What we release}
\label{sec:release}


Everything needed to reproduce this paper is public: the benchmark, the
evaluation code, all 5,194 model responses, every judge transcript, the
preference dataset, and the post-trained model. Appendix~\ref{app:artifacts}
gives the locations and the licences. We do not redistribute the embryo images,
which remain with the source dataset \citep{kromp2023} under CC BY 4.0; our
release rebuilds every case from that source deterministically.

The study cost \$813 as recorded in those artifacts: \$432 for 51.5 hours of a
two-H100 instance covering preference-data preparation, the hyperparameter
search, training, and local inference, plus \$328 of API spend on the test split
and \$53 on the held-out split. A further \$15 was spent re-running judge calls
that the two defects in Section~\ref{sec:setup} had invalidated; those superseded
calls are not counted in the per-model figures, which reflect the final scores
only.


\bibliography{references}

\appendix
\section{Case prompt template}
\label{app:prompt}
Every model receives this prompt together with the embryo image. The example is case IVF-BENCH-0001; laboratory and patient values change per case, the structure does not.

\begin{lstlisting}
You are an expert reproductive medicine specialist reviewing an IVF case.

## Embryo Image
[IMAGE ATTACHED]

## Laboratory Data
- Gardner Score: 1DD (Expansion: blastocyst, blastocoel at least half the embryo, ICM: degenerative or absent, TE: degenerative or absent)
- COCs retrieved: 8
- MII oocytes: 6
- Transfer day: Day 4

## Patient Profile
- Female age: 39 | BMI: 28.1
- AMH: 2.0 ng/mL | FSH: 5.8 IU/L
- Diagnosis: Tubal Factor
- Stimulation protocol: Long Agonist
- Previous IVF cycles: 3 (no_pregnancy, clinical_miscarriage, biochemical_pregnancy)
- Partner age: 37
- Sperm: 88.2 M/mL, 78.2% motile
- Endometrial thickness: 7.0 mm
- Lifestyle: never smoker; moderate drinker (4-6/wk); very active (daily exercise); standard diet; moderate stress

## Task
Provide a comprehensive clinical assessment including:
1. Morphological analysis of the blastocyst image
2. Integration of patient factors with embryo quality
3. Risk factors and favorable indicators for this case
4. Estimated implantation probability range with reasoning
5. Recommendations for this transfer and future cycle planning
\end{lstlisting}


\section{Rubric definitions}
\label{app:rubrics}
The judge returns a score and a one-line justification for each rubric, plus the numeric probability it extracted from the response. Rubric text The rubric text below is what the judge sees, verbatim.

\paragraph{Morphological Accuracy}
How accurately does the response describe the embryo's morphological features as visible in the image? Consider: expansion stage assessment, ICM description, TE description, zona pellucida observations, overall blastocyst appearance.
\begin{itemize}
\item \textbf{1}: Major errors or fabricated features not visible in the image
\item \textbf{2}: Several inaccuracies in morphological description
\item \textbf{3}: Mostly accurate with some minor errors or omissions
\item \textbf{4}: Accurate description with good detail
\item \textbf{5}: Precise, thorough morphological analysis matching the image
\end{itemize}

\paragraph{Clinical Integration}
How well does the response integrate patient clinical factors (age, AMH, BMI, diagnosis, previous outcomes, sperm parameters) with the embryo assessment? Does it consider how these factors interact with embryo quality?
\begin{itemize}
\item \textbf{1}: No integration of clinical factors with embryo assessment
\item \textbf{2}: Mentions clinical factors but doesn't connect them to embryo quality
\item \textbf{3}: Some integration but misses key interactions
\item \textbf{4}: Good integration of most relevant clinical factors
\item \textbf{5}: Comprehensive integration showing how clinical factors affect prognosis
\end{itemize}

\paragraph{Reasoning Coherence}
Is the clinical reasoning logical, consistent, and free of contradictions? Does the probability estimate follow from the evidence presented? Are conclusions supported by the observations?
\begin{itemize}
\item \textbf{1}: Contradictory or illogical reasoning
\item \textbf{2}: Significant logical gaps or inconsistencies
\item \textbf{3}: Generally coherent with minor logical issues
\item \textbf{4}: Well-structured reasoning with clear logic
\item \textbf{5}: Exemplary reasoning chain with all conclusions well-supported
\end{itemize}

\paragraph{Guideline Alignment}
Does the response align with established IVF clinical guidelines (ESHRE, ASRM, Gardner grading criteria)? Are the recommendations consistent with evidence-based reproductive medicine?
\begin{itemize}
\item \textbf{1}: Recommendations contradict established guidelines
\item \textbf{2}: Several recommendations not aligned with guidelines
\item \textbf{3}: Mostly aligned with some deviations
\item \textbf{4}: Well-aligned with current clinical guidelines
\item \textbf{5}: Fully aligned, demonstrates deep knowledge of current guidelines
\end{itemize}

\paragraph{Actionability}
Are the recommendations specific, practical, and useful for a clinician? Does the response provide clear next steps for this transfer and future cycle planning? Are recommendations appropriately nuanced?
\begin{itemize}
\item \textbf{1}: No actionable recommendations or only generic advice
\item \textbf{2}: Vague recommendations with little practical value
\item \textbf{3}: Some useful recommendations but lacking specificity
\item \textbf{4}: Clear, specific recommendations for this case
\item \textbf{5}: Highly specific, nuanced recommendations with alternatives considered
\end{itemize}


\section{Generated patient-context distributions}
\label{app:distributions}
These are the fields that no public dataset carries alongside embryo images and outcomes, and that we therefore generate. Each is drawn independently, from the distribution shown, seeded deterministically from the case identifier. Sources are the nearest published population estimate we could find for an IVF-treated cohort; where registries disagree by geography we chose a Western-representative central value.

\begin{table}[H]
\centering\footnotesize
\begin{tabular}{p{3.1cm}p{7.4cm}p{2.6cm}}
\toprule
\textbf{Field} & \textbf{Distribution} & \textbf{Source} \\
\midrule
BMI (kg/m$^2$) & Truncated normal, mean 25.0, SD 5.0, range [17, 45] & \citet{eshre2025} \\
Basal FSH (IU/L) & Truncated normal, mean 6.0 at age 30, $+0.15$ per year & \citet{eshre2025} \\
Diagnosis & Categorical: tubal 0.25, male 0.25, unexplained 0.15, PCOS 0.12, endometriosis 0.10, DOR 0.05, uterine 0.03, other 0.05 & \citet{eshre2025} \\
Stimulation protocol & Categorical: antagonist 0.70, long agonist 0.20, short agonist 0.05, natural 0.03, other 0.02 & \citet{yoo2026} \\
Previous cycles & Categorical: 0 at 0.40, then 0.25, 0.15, 0.10, 0.05, 0.05 & \citet{vanderborght2025} \\
Previous outcomes & Categorical per cycle: no pregnancy 0.45, biochemical 0.20, miscarriage 0.15, live birth 0.15, ectopic 0.05 & \citet{vanderborght2025} \\
Partner age (years) & Female age $+$ Normal(2.4, 4.3), clipped to [20, 65] & \citet{jelinkova2026} \\
Sperm concentration (M/mL) & Truncated normal, mean 50.0, SD 25.0, range [5, 200] & \citet{who2021semen} \\
Sperm motility (\%) & Truncated normal, mean 55.0, SD 15.0, range [10, 95] & \citet{who2021semen} \\
Smoking & Categorical: never 0.65, former 0.25, current 0.10 & \citet{dodge2015} \\
Alcohol & Categorical: none 0.40, light 0.35, moderate 0.20, heavy 0.05 & \citet{dodge2015} \\
Physical activity & Categorical: sedentary 0.30, light 0.25, moderate 0.25, active 0.15, very active 0.05 & \citet{sherwin2022} \\
Diet & Categorical: standard 0.50, health-conscious 0.25, Mediterranean 0.15, vegetarian 0.05, restricted 0.05 & \citet{gaskins2019} \\
\bottomrule
\end{tabular}
\caption{Generated patient-context fields. Because each is sampled independently of the recorded outcome and of the other fields, they carry no outcome signal and no inter-field correlation.}
\end{table}


\section{Code and data availability}
\label{app:artifacts}
The benchmark, the evaluation code, every model response, every judge transcript, and the analysis that produced each interval reported here are available at \url{https://github.com/The-Fertility-Plan/ivf-bench}. The post-trained model is at \url{https://huggingface.co/thefertilityplan/ivf-bench-qwen9b-vlm-orpo} and the preference dataset at \url{https://huggingface.co/datasets/thefertilityplan/ivf-bench-orpo-qwen9b-clipped}. Code is released under Apache 2.0.

We do not redistribute the embryo images. They belong to the Kromp blastocyst dataset \citep{kromp2023}, are available under CC BY 4.0 at \url{https://doi.org/10.6084/m9.figshare.20123153.v3}, and our release rebuilds every case from that source deterministically.

To cite this work:

\begin{lstlisting}
@article{correa2026ivfbench,
  author  = {Correa, Andrew G. A. and Yoon, Brittany},
  title   = {{IVF-Bench}: Vision-Language Models Explain Embryo Cases Well and
             Predict Outcomes Poorly},
  journal = {arXiv preprint},
  year    = {2026}
}
\end{lstlisting}


\end{document}
